\section{Observabilidad y evaluación ISO/IEC 25010}
\label{sec:observabilidad-iso}

La Entrega 4 incorpora observabilidad como una capacidad transversal del
sistema distribuido y no como sustituto de las pruebas. Métricas, trazas y logs
responden preguntas distintas: las métricas resumen comportamiento agregado,
las trazas reconstruyen el recorrido de una solicitud y los logs conservan
eventos discretos con contexto. Esta separación coincide con las señales
definidas por OpenTelemetry \cite{opentelemetry2026signals}; combinarlas mejora
el diagnóstico, pero la sola presencia de herramientas no certifica calidad ni
garantiza correlación perfecta.

\subsection{Arquitectura de observabilidad}
\label{subsec:arquitectura-observabilidad}

Los cinco servicios Spring exponen Actuator y Prometheus, y configuran
instrumentación y exportación OTLP. El flujo implementado es el siguiente:

\begin{center}
\begin{verbatim}
servicios -- metricas Prometheus ----------> Prometheus --\
         \-- trazas OTLP --> Collector --> Tempo ----------> Grafana
          \- logs OTLP ----> Collector --> Loki -----------/
cAdvisor -- recursos de contenedores ------> Prometheus ---/
\end{verbatim}
\end{center}

Prometheus consulta cada quince segundos los cinco endpoints
\texttt{/actuator/prometheus}, cAdvisor, el Collector y los tres nodos
CockroachDB. Grafana dispone de fuentes provisionadas para Prometheus, Tempo y
Loki. Tempo conserva trazas localmente durante 24 horas; Loki almacena logs, y
cAdvisor aporta CPU y memoria de contenedores. El Collector recibe OTLP por
gRPC 4317 y HTTP 4318, aplica el procesador \texttt{batch} y separa tres
pipelines: trazas hacia Tempo, métricas hacia un exporter Prometheus en 8889 y
logs hacia Loki. OpenTelemetry funciona aquí como estándar, instrumentación y
canal de transporte; los backends y Grafana proporcionan almacenamiento y
consulta, pues OpenTelemetry no es por sí mismo un APM completo.

\subsubsection{Métricas Prometheus}

Micrometer produce la familia
\texttt{http\_server\_requests\_seconds}: el contador permite calcular volumen y
tasa por estado, y los buckets habilitan percentiles con
\texttt{histogram\_quantile}. En Reservas una regla de grabación crea el alias
\texttt{http\_request\_duration\_seconds} exigido por el protocolo a partir de
la familia estándar. Usuarios y Académico agregan además
\texttt{http\_requests\_total} con ruta, método y estado; Auth registra éxitos,
fallos, refresh y logout; Reservas publica eventos de negocio de baja
cardinalidad después del commit transaccional. Las series de CockroachDB y
cAdvisor complementan la vista con estado de nodos, rangos no disponibles y
consumo de recursos.

La tasa se obtiene con \texttt{rate} sobre contadores y los 5xx mediante el
filtro \texttt{status=\textasciitilde{}"5.."}. Para latencia, los buckets se
agregan por \texttt{le} antes de calcular p50, p95 o p99. Un percentil de
histograma es una estimación dependiente de los buckets y de la ventana elegida;
no equivale necesariamente al percentil calculado por el cliente de carga.

\subsubsection{Trazas y logs estructurados}

La configuración de cada servicio asigna \texttt{OTEL\_SERVICE\_NAME}, exporta
trazas y logs por OTLP y habilita muestreo de trazas. Los appenders Logback se
instalan sobre el SDK autoconfigurado y el formato de consola es JSON tipo
Logstash. La configuración añade el servicio y renombra el contexto activo como
\texttt{trace\_id} y \texttt{span\_id}; el encoder incorpora también timestamp,
nivel, logger y mensaje. Tempo puede enlazar una traza con Loki usando ambos
identificadores, aunque un log emitido fuera de un span no necesariamente los
contiene. Por esa razón se habla de correlación disponible, no universal.

Los registros no deben transportar credenciales. En particular, el callback de
FCM informa que recibió un token nuevo, pero no escribe su valor. Esta decisión
preserva utilidad operativa sin convertir el sistema de logs en un canal de
exposición. Los eventos estructurados y ricos en contexto facilitan consultas y
diagnóstico de sistemas distribuidos \cite{majors2022observability}, pero su
calidad depende de atributos consistentes y de la propagación efectiva.

\subsection{Dashboard consolidado}
\label{subsec:dashboard-grafana}

\texttt{ops/grafana/pfc-dashboard.json} contiene exactamente seis paneles. Las
consultas PromQL son sintácticamente coherentes con las familias configuradas;
la Tabla~\ref{tab:paneles-grafana} distingue lo que mide cada una y evita
atribuir al panel una semántica que no posee.

\begin{table}[ht]
  \centering
  \small
  \caption{Paneles reales del dashboard consolidado.}
  \label{tab:paneles-grafana}
  \begin{tabular}{|p{3.0cm}|p{2.0cm}|p{3.6cm}|p{3.8cm}|}
    \hline
    \textbf{Panel} & \textbf{Fuente} & \textbf{Métrica/consulta} & \textbf{Propósito} \\
    \hline
    Request rate & Prometheus & \texttt{rate(http\_server\_requests\_seconds\_count)} & Comparar volumen por job en cinco minutos. \\
    \hline
    Latencia p50/p95/p99 & Prometheus & \texttt{histogram\_quantile} sobre buckets HTTP & Observar distribución y cola de latencia. \\
    \hline
    Errores HTTP 5xx & Prometheus & Tasa del contador HTTP con estado 5xx & Detectar fallos de servidor por job. \\
    \hline
    CockroachDB & Prometheus & \texttt{liveness\_livenodes}; \texttt{ranges\_unavailable} & Vigilar nodos vivos y rangos no disponibles. \\
    \hline
    Recursos de contenedores & Prometheus/cAdvisor & CPU y memoria por identificador de contenedor & Identificar presión de recursos. \\
    \hline
    Latencia E2E móvil & Prometheus & \texttt{histogram\_quantile} p95 sobre \texttt{app\_mobile\_http\_latency\_seconds\_bucket} & Medir latencia HTTP observada desde Android, reportada a Reservas y expuesta por Prometheus. \\
    \hline
  \end{tabular}
\end{table}

El sexto panel ya no usa el Gateway como proxy. La aplicación Android mide el
tiempo HTTP observado por OkHttp, lo reporta a
\texttt{POST /api/v1/observabilidad/mobile/http-latency} y
Reservas expone las familias \texttt{app.mobile.http.latency} y
\texttt{app.mobile.http.requests} para Prometheus. El panel calcula p95 con
\texttt{app\_mobile\_http\_latency\_seconds\_bucket} y el job
\texttt{reservas-solicitudes-service}. La evidencia runtime final sigue
pendiente: aún debe capturarse el dashboard en la VM con esta serie poblada bajo
carga.

La Figura~\ref{fig:panel-monitoreo-carga-real} muestra el dashboard consolidado
versionado como evidencia visual de su estructura y provisionamiento. La imagen
actual no se atribuye a una nueva validación de 50 usuarios, cinco minutos ni
22,64 solicitudes por segundo, porque esa relación debe provenir de la evidencia
runtime correspondiente. Cuando se ejecute la validación final en la VM será
necesaria una nueva captura bajo carga que muestre la serie
\texttt{app\_mobile\_http\_latency\_seconds\_bucket} poblada. Esta captura no
sustituye ni modifica las diez repeticiones nominales analizadas en la
Sección~\ref{subsec:eficiencia-nominal}.

\begin{figure*}[htbp]
  \centering
  \includegraphics[width=\textwidth]{panel-monitoreo.png}
  \caption{Dashboard de observabilidad de SCLI versionado como evidencia visual
  de estructura y provisionamiento.}
  \label{fig:panel-monitoreo-carga-real}
\end{figure*}

\subsection{Marco y protocolo ISO/IEC 25010}
\label{subsec:marco-iso25010}

La evaluación usa el modelo de calidad de producto ISO/IEC 25010:2023
\cite{iso25010_2023}, pero no pretende cubrir todas sus características. Se
seleccionaron eficiencia de desempeño, fiabilidad, seguridad, mantenibilidad y
compatibilidad. La disponibilidad se trata como parte de fiabilidad y exige
observación temporal propia; no se infiere de una tasa puntual de errores.

El protocolo nominal define diez repeticiones equivalentes con 50 usuarios. La
eficiencia dura cinco minutos y exige p95 menor que 500 ms. La fiabilidad
preveía una hora por repetición y menos de 1\% de HTTP 5xx. La regla se fijó
antes de observar resultados: r1 y r10 se conservan, mientras r2--r9 forman las
ocho muestras centrales. Esto trata de forma uniforme los bordes del lote y,
sobre todo, impide elegir exclusiones después de conocer los valores; no se
eliminaron ejecuciones desfavorables por conveniencia.

Para cada métrica central se usa $n=8$, desviación estándar muestral con siete
grados de libertad e intervalo bilateral
$\bar{x}\pm t_{0,975;7}s/\sqrt{8}$, con $t=2,364624251$. La decisión principal
es conservadora: el límite superior del IC95 debe ser menor que el umbral. Esta
estructura y la conservación de resultados fallidos son coherentes con un
proceso de prueba documentado y repetible \cite{iso29119_2_2021}.

\subsection{Preguntas de investigación y respuestas}
\label{subsec:preguntas-investigacion}

Los criterios, umbrales y diseño experimental fueron declarados antes de la
ejecución. Las preguntas PI1/PI2 se formalizaron posteriormente en el manuscrito
utilizando esos criterios preexistentes. El historial de Git sitúa el protocolo
inicial y los umbrales de p95 y HTTP 5xx en el commit
\texttt{5b1c498} del 18 de agosto de 2026. El umbral secundario de p99 fue
incorporado al analizador en \texttt{e1c82c7} el 24 de agosto, antes de que la
campaña de eficiencia se ejecutara el 30 de agosto. En ambos casos la regla
permite una respuesta negativa: se exige que el límite superior del IC95 sea
estrictamente menor que el umbral.

\textbf{PI1.} ¿Al fijar como variable manipulada una carga nominal de 50
usuarios concurrentes durante cinco minutos sobre las consultas autenticadas
de solo lectura de Reservas/Solicitudes, los límites superiores de los IC95 de
la latencia p95 y p99 observada por Locust permanecen por debajo de los
umbrales previamente declarados de 500 ms y 750 ms, respectivamente, en las
ocho repeticiones centrales r2--r9 del despliegue SCLI evaluado?

\textbf{Respuesta a PI1.} Sí, para el escenario observado en la campaña
poblada post-evaluación, que constituye el cierre oficial de PI1. Se midieron
50 usuarios durante cinco minutos en diez repeticiones independientes, con un
dataset controlado de 20 reservas y ambas rutas de negocio:
\texttt{GET /api/v1/reservas} y \texttt{GET /api/v1/reservas/\{id\}}.
Las ocho repeticiones centrales r2--r9 contienen 59\,278 GET:
44\,546 listados y 14\,732 consultas por id, con 0 respuestas HTTP 401 y
0 respuestas HTTP 5xx. Para p95 se obtiene media 9,952465 ms, desviación
muestral 0,426872 ms e IC95 [9,595591; 10,309339] ms; para p99, media
19,744969 ms, desviación muestral 3,437219 ms e IC95
[16,871382; 22,618556] ms. Los límites superiores satisfacen
$10,309339 < 500$ ms y $22,618556 < 750$ ms, respectivamente, según la
regla t de ocho muestras del diseño existente.

El software desplegado fue
\texttt{dd2e1923d7635857f4020c7ef5c8b6efee363d7a}. El analizador reconstruye
los percentiles por repetición desde todos los eventos GET de negocio de
\texttt{locust\_requests.csv}, sin filtrar por código HTTP, latencia ni
éxito/fallo; login y refresh quedan fuera por pertenecer a Auth.
Las fuentes oficiales son
\path{experimentos/resultados/iso25010-eficiencia-poblada.csv} y
\path{experimentos/resultados/raw/eficiencia_nominal_50u_5m_poblada/};
el protocolo está en
\path{experimentos/protocolo-e4-eficiencia-poblada.md}.
La conclusión queda acotada al escenario, datos, carga y entorno medidos:
20 reservas no implican representatividad universal ni reconstrucción exacta
del dataset histórico. La campaña vacía del 30 de agosto se conserva como
antecedente en la Sección~\ref{subsec:eficiencia-nominal}, sin constituir la
respuesta definitiva a toda la población declarada.

\textbf{PI2.} ¿Al fijar como variable manipulada una carga nominal de 50
usuarios concurrentes durante una hora sobre las consultas autenticadas de
solo lectura de Reservas/Solicitudes, el límite superior del IC95 de la tasa
observada de respuestas HTTP 5xx permanece por debajo del umbral previamente
declarado de 1\,\% en las ocho repeticiones centrales r2--r9 del despliegue
SCLI evaluado, y qué comportamiento operacional adicional registra Locust?

\textbf{Respuesta a PI2.} Sí para el indicador formal acotado en la
campaña correctiva poblada post-evaluación. Las diez repeticiones sumaron
895\,136 GET de negocio, de los cuales 224\,164 fueron consultas por id, con
0 respuestas HTTP 401 y 42 respuestas HTTP 5xx. Los 5xx observados se
preservaron: r7 registró 1, r8 registró 21 y r9 registró 20.

Sobre r2--r9 se analizaron 716\,099 GET de negocio. La tasa 5xx tuvo media
0,005870\,\%, desviación muestral 0,010535\,\% e IC95
[-0,002938; 0,014677]\,\%. El intervalo t se reporta sin truncamiento;
su límite superior, 0,014677\,\%, permanece por debajo del 1\,\%
prerregistrado.

Como métricas descriptivas, p95 tuvo media 11,256521 ms e IC95
[11,069649; 11,443394] ms, y p99 media 23,945109 ms e IC95
[22,605247; 25,284971] ms. Estos percentiles no constituyen criterios
adicionales de aceptación E2.

La campaña histórica con 668\,367 HTTP 401 permanece publicada como
antecedente del defecto de renovación. Una primera campaña correctiva con
refresh, ejecutada sobre \texttt{b94b7af7ebab2510c16eae0c70593664763de1e7},
también permanece preservada, pero no se utiliza como cierre de PI2 porque su
dataset de reservas estaba vacío y no ejercitó
\texttt{GET /api/v1/reservas/\{id\}}.

La campaña nueva usa una población controlada de 20 reservas. No se afirma que
sea una reconstrucción exacta del dataset histórico, por lo que la
comparabilidad estricta entre datasets queda limitada y declarada.

La tasa 5xx no se interpreta como disponibilidad temporal. El consolidado
oficial es
\path{experimentos/resultados/iso25010-correctiva-poblada.csv};
la evidencia raw está en
\path{experimentos/resultados/raw/fiabilidad_nominal_50u_1h_refresh_poblada/}.
El analizador reconstruye las métricas directamente desde
\texttt{locust\_requests.csv}.

\subsubsection{Proveniencia e integridad}

Las campañas deben distinguirse por su procedencia. La campaña histórica de
fiabilidad fue ejecutada sobre
\texttt{061a1050a94e1bd30d81b30c47c7e818005a33bb}; la primera campaña con
refresh y dataset vacío se ejecutó sobre
\texttt{b94b7af7ebab2510c16eae0c70593664763de1e7}. Ambas se conservan como
antecedentes.

La campaña que cierra la corrección de fiabilidad es
\path{experimentos/resultados/raw/fiabilidad_nominal_50u_1h_refresh_poblada/}
y midió el software desplegado
\texttt{dd2e1923d7635857f4020c7ef5c8b6efee363d7a}. Se ejecutó con Locust
2.31.6 y Python 3.12.3 en la VM \texttt{pfc-pruebas-e1}.

Cada repetición poblada conserva un manifiesto SHA-256. El manifiesto global
\path{experimentos/resultados/SHA256SUMS} contiene 1.339 entradas.
Incluye los 550 archivos actualmente versionados de la campaña correctiva
anterior, su consolidado
\path{experimentos/resultados/iso25010-correctiva.csv}, las 319
evidencias de la campaña poblada y el consolidado
\path{experimentos/resultados/iso25010-correctiva-poblada.csv}.

La población de 20 reservas y su hash están fijados en la evidencia de
preflight. No se atribuye a esa población identidad con el dataset histórico;
la desviación se declara como amenaza a la comparabilidad.

\subsection{Eficiencia nominal: 50 usuarios durante cinco minutos}
\label{subsec:eficiencia-nominal}

\subsubsection{Cierre post-evaluación poblado}

La respuesta oficial a PI1 de la
Sección~\ref{subsec:preguntas-investigacion} utiliza la campaña poblada:
diez repeticiones con 74\,098 GET, incluidos 55\,800 listados y 18\,298
consultas por id, con 0 HTTP 401 y 0 HTTP 5xx. Sus estadísticos centrales y
la comparación con ambos umbrales se presentan allí para evitar duplicación.
El dataset controlado de 20 reservas conserva el SHA-256
\texttt{f2c07cc4698715939a9eaf6b47e694ccdf39fcd2de0502aec6164a90321f9b78}.

\subsubsection{Antecedente histórico: campaña del 30 de agosto}

La campaña original utilizó un dataset vacío y no ejercitó
\texttt{GET /api/v1/reservas/\{id\}}. Esta amenaza de validez de constructo
impide generalizar sus resultados a toda la población declarada en PI1.
Se conservan sus cifras y trazabilidad en
\path{experimentos/resultados/raw/eficiencia_nominal_50u_5m/}.
Las diez repeticiones históricas terminaron con código cero y registraron 0 respuestas HTTP
5xx. La Tabla~\ref{tab:repeticiones-eficiencia} transcribe el CSV verificado; r1
y r10 se muestran porque forman parte de la evidencia, aunque no intervienen en
el análisis central.

\begin{table}[ht]
  \centering
  \small
  \caption{Antecedente histórico: eficiencia nominal con dataset vacío medida por Locust.}
  \label{tab:repeticiones-eficiencia}
  \begin{tabular}{|c|r|r|r|r|c|}
    \hline
    \textbf{Rep.} & \textbf{Requests} & \textbf{HTTP 5xx} & \textbf{p95 (ms)} & \textbf{p99 (ms)} & \textbf{Uso} \\
    \hline
    r1  & 6991 & 0 & 87 & 1000 & Conservada \\
    r2  & 7180 & 0 & 24 & 40   & Central \\
    r3  & 7101 & 0 & 80 & 1000 & Central \\
    r4  & 7162 & 0 & 24 & 45   & Central \\
    r5  & 7155 & 0 & 54 & 740  & Central \\
    r6  & 7209 & 0 & 31 & 85   & Central \\
    r7  & 7164 & 0 & 58 & 410  & Central \\
    r8  & 7159 & 0 & 27 & 110  & Central \\
    r9  & 7111 & 0 & 66 & 540  & Central \\
    r10 & 7231 & 0 & 23 & 88   & Conservada \\
    \hline
  \end{tabular}
\end{table}

Las ocho repeticiones centrales históricas suman 57\,241 consultas GET de
listado. Para p95, los valores 24, 80, 24, 54, 31, 58, 27 y 66 ms producen media
45,500000 ms, desviación muestral 21,764978 ms e IC95
[27,304023; 63,695977] ms. Como el límite superior 63,695977 ms es menor que
500 ms, satisface el umbral numérico en el escenario histórico vacío, sin
resolver la cobertura de ambas rutas requerida por PI1.

El p99 se analiza como indicador secundario. Los valores 40, 1000, 45, 740, 85,
410, 110 y 540 ms producen media 371,250000 ms, desviación 364,032475 ms e IC95
[66,911235; 675,588765] ms. El límite superior permanece bajo el umbral
secundario de 750 ms, pero la dispersión es elevada y la distancia de
74,411235 ms al umbral no se interpreta como margen holgado. Este resultado
histórico no constituye el cierre definitivo de PI1; sus raws y umbrales se
conservan intactos.

Para trazabilidad, el análisis histórico de la fila \texttt{Aggregated} mezcló
57\,241 GET con 400 login en r2--r9. Produjo p95 media 57,500000 ms e IC95
[26,456064; 88,543936] ms, y p99 media 624,000000 ms e IC95
[116,881810; 1131,118190] ms; p99 no cumplía bajo esa población incorrecta.

Las ocho ventanas centrales contienen 0 HTTP 5xx de 57\,241 consultas GET, es
decir, 0\%, frente al objetivo menor que 1\%. Esto demuestra ausencia de 5xx en
el escenario observado; no demuestra disponibilidad anual, ausencia de otros
errores ni fiabilidad absoluta.

\subsubsection{Contraste histórico con Prometheus}

Las series Prometheus de r2--r9 contienen p95 de 10,004445; 30,067916;
9,949600; 21,021640; 13,281961; 22,143963; 11,358052 y 28,439426 ms. Su media es
18,283375 ms, la desviación 8,240924 ms y el IC95 [11,393791; 25,172960] ms;
las diez ventanas registran cero 5xx. Locust mide extremo a extremo desde el
generador, incluyendo red, Gateway y serialización. Prometheus mide el histograma
interno de Reservas para su ventana. Son puntos de observación complementarios,
no resultados intercambiables ni fuentes que deban forzarse a coincidir.

\subsection{Rampa exploratoria de 0 a 200 usuarios}
\label{subsec:rampa-exploratoria}

La rampa de diez minutos produjo 29,217 solicitudes, una respuesta HTTP 500 en
\texttt{GET /api/v1/reservas} y una tasa exacta de 0,00342266\%. Locust registró
p95 de 180 ms, p99 de 1100 ms y máximo de 14024,163663 ms. El proceso terminó
con estado fallido porque se configuró para devolver error ante cualquier fallo,
no porque la mayoría de solicitudes hubiera fallado.

La tasa 5xx permaneció por debajo de 1\% y el p95 fue moderado mientras la carga
crecía, pero p99 y el máximo confirman una cola larga y picos aislados. Esta
prueba complementa el estudio de capacidad progresiva; no pertenece al CSV
central ni sustituye el escenario nominal ISO.

\subsection{Fiabilidad nominal y alcance de la disponibilidad}
\label{subsec:fiabilidad-disponibilidad}

La evolución experimental comprende tres campañas. La campaña histórica
\texttt{061a1050...} conserva 668\,367 HTTP 401 causados por expiración de
sesión. La primera campaña con refresh, \texttt{b94b7af...}, eliminó ese
defecto de sesión pero se ejecutó con una tabla de reservas vacía, de modo que
no ejercitó la consulta por id y se conserva únicamente como antecedente.

La campaña correctiva poblada completó diez repeticiones de aproximadamente una
hora con 50 usuarios y \texttt{spawn-rate} 10 usuarios/s sobre el software
desplegado
\texttt{dd2e1923d7635857f4020c7ef5c8b6efee363d7a}. El análisis central
utiliza r2--r9.

Se fijó una población controlada de 20 reservas. Esta población corrige la
degeneración vacía, pero no se declara como reconstrucción exacta del dataset
histórico.

Las diez repeticiones contienen 895\,136 GET de negocio: 670\,972 de listado
y 224\,164 por id. Registraron 0 HTTP 401 y 42 HTTP 5xx. Los errores 5xx no
se descartaron: r7 conservó 1, r8 conservó 21 y r9 conservó 20.

Además, todas las repeticiones mantuvieron consultas por id después de 900 s,
lo que aporta evidencia de carga autenticada sostenida más allá del TTL que
había degradado la campaña histórica.

Sobre r2--r9 se analizaron 716\,099 GET de negocio. La tasa HTTP 5xx tuvo media
0,005870\,\%, desviación muestral 0,010535\,\% e IC95
[-0,002938; 0,014677]\,\%. El intervalo Student t no está acotado a valores
físicamente posibles, por lo que el límite inferior negativo se reporta sin
truncarlo. El límite superior 0,014677\,\% es menor que 1\,\%; por tanto,
\textbf{E2 cumple} el criterio acotado prerregistrado en la campaña poblada.

Como métricas descriptivas, p95 presentó media 11,256521 ms, desviación
muestral 0,223526 ms e IC95 [11,069649; 11,443394] ms. P99 presentó media
23,945109 ms, desviación muestral 1,602666 ms e IC95
[22,605247; 25,284971] ms. Estos percentiles no constituyen criterios
adicionales de aceptación E2.

Las métricas se reconstruyen desde \texttt{locust\_requests.csv} sobre todos
los GET de negocio, sin excluir observaciones por código HTTP, éxito, fallo o
latencia. Login y refresh se contabilizan por separado.

El código de salida de Locust se conserva como evidencia operacional, pero no
es por sí solo criterio de exclusión cuando refleja fallos HTTP observados.
Así, r7--r9 permanecen en la población aun cuando finalizaron con exit code 1;
excluirlas habría eliminado observaciones desfavorables.

La tasa HTTP 5xx y la disponibilidad temporal no son equivalentes. Para afirmar
una disponibilidad de al menos 99,5\,\% sería necesario operacionalizar tiempo
apto frente a tiempo total. En consecuencia, E2 es concluyente para su criterio
acotado de errores, mientras la disponibilidad temporal permanece
\textbf{no concluyente}.

\subsection{Experimento de arbitraje concurrente S0--S4}
\label{subsec:arbiter-resultados}

El experimento ARBITER comparó cinco comportamientos seleccionables sin alterar
el camino productivo predeterminado: S0, escritura sin arbitraje usada como
control negativo; S1, control optimista; S2, control pesimista; S3, coordinador
de aplicación elegido mediante el algoritmo Bully y orden lógico de Lamport; y
S4, transacción \texttt{SERIALIZABLE} en CockroachDB respaldada por el consenso
del clúster. S3 no utiliza Raft: su líder pertenece a la aplicación. En S4 no
existe un coordinador Bully; la coordinación corresponde a la persistencia
distribuida del clúster.

La campaña del SHA
\texttt{e43967eda31410b7fef060cfc7903cea94efbe96} planificó 130 corridas. Se
completaron 123, siete fallaron y ninguna fue clasificada como inválida. La regla
prerregistrada excluye las repeticiones 1 y 10 de cada configuración, lo que
produce 104 posiciones analíticas teóricas. Cinco posiciones centrales de
Esc-1 también fallaron; por ello el conjunto realmente analizable contiene 99
corridas completadas. La Tabla~\ref{tab:arbiter-conteos} evita confundir esas
cuatro cantidades.

\begin{table}[ht]
  \centering
  \small
  \caption{Conteos auditados de la campaña ARBITER.}
  \label{tab:arbiter-conteos}
  \begin{tabular}{|l|r|r|r|r|}
    \hline
    \textbf{Escenario} & \textbf{Planificadas} & \textbf{Completadas} & \textbf{Fallidas} & \textbf{Analíticas válidas} \\
    \hline
    Esc-1, nominal & 10 & 3 & 7 & 3 \\
    \hline
    Esc-2, contención media & 50 & 50 & 0 & 40 \\
    \hline
    Esc-3, contención máxima & 50 & 50 & 0 & 40 \\
    \hline
    Esc-4, caída controlada & 20 & 20 & 0 & 16 \\
    \hline
    \textbf{Total} & \textbf{130} & \textbf{123} & \textbf{7} & \textbf{99} \\
    \hline
  \end{tabular}
\end{table}

S0 exhibió dobles adjudicaciones en todas las muestras analíticas de
contención: 8/8 en Esc-2 y 8/8 en Esc-3. Para S1--S4 no se observaron dobles
adjudicaciones en las corridas ejecutadas; este resultado acotado no constituye
una garantía de aislamiento absoluto. La Tabla~\ref{tab:arbiter-estrategias}
resume las 96 muestras analíticas de contención y caída, separadas de las tres
muestras nominales válidas de Esc-1.

\begin{table}[ht]
  \centering
  \small
  \caption{Resultado principal por estrategia en muestras analíticas.}
  \label{tab:arbiter-estrategias}
  \begin{tabular}{|c|r|l|l|}
    \hline
    \textbf{Estrategia} & \textbf{$n$} & \textbf{Dobles adjudicaciones} & \textbf{Interpretación} \\
    \hline
    S0 & 16 & Observadas en 16/16 & Control negativo efectivo \\
    \hline
    S1 & 16 & No observadas & Resultado limitado a las corridas \\
    \hline
    S2 & 16 & No observadas & Resultado limitado a las corridas \\
    \hline
    S3 & 24 & No observadas & Incluye 8 corridas de caída \\
    \hline
    S4 & 24 & No observadas & Incluye 8 corridas de caída \\
    \hline
  \end{tabular}
\end{table}

La Tabla~\ref{tab:arbiter-http} presenta el censo HTTP de la población central
r2--r9, incluyendo las cinco corridas nominales fallidas de ese intervalo.
Para Esc-2, Esc-3 y Esc-4 se deriva de los registros \texttt{REQUEST} de los
JSON raw: el código explícito se conserva en
\texttt{allocation.backend\_response.httpStatus}; una respuesta normal del
endpoint se registra como 200. El estado de negocio \texttt{REJECTED} no
equivale a un error HTTP. En Esc-1 los JSON identifican las corridas, pero
contienen \texttt{records} vacío: sus conteos se reconstruyen desde
\texttt{locust\_stats.csv} y \texttt{locust\_failures.csv} en los directorios
\texttt{esc1-production-r02} a \texttt{esc1-production-r09}. Se incluye solo
\texttt{GET /api/v1/reservas}, excluyendo los 400 login de Auth.
El harness exige HTTP 200; al coincidir el detalle de fallos con
\texttt{Failure Count}, los GET restantes se contabilizan como 200.
En esta población central delimitada, todos los errores HTTP observados fueron
500; no se encontraron 4xx ni otros códigos. La tasa de error es el cociente
entre solicitudes no 2xx y solicitudes totales, no una media de tasas por corrida.

\begin{table}[ht]
  \centering
  \scriptsize
  \caption{Censo HTTP ARBITER en la población central r2--r9.}
  \label{tab:arbiter-http}
  \begin{tabular}{|c|c|r|r|r|r|r|}
    \hline
    \textbf{Esc.} & \textbf{Estrategia} & \textbf{Requests} & \textbf{2xx} & \textbf{4xx} & \textbf{5xx} & \textbf{Error} \\
    \hline
    1 & Producción & 58\,040 & 57\,969 & 0 & 71 & 0,1223\% \\
    \hline
    2 & S0 & 400 & 400 & 0 & 0 & 0\% \\
    2 & S1 & 400 & 400 & 0 & 0 & 0\% \\
    2 & S2 & 400 & 348 & 0 & 52 & 13,0000\% \\
    2 & S3 & 400 & 400 & 0 & 0 & 0\% \\
    2 & S4 & 400 & 396 & 0 & 4 & 1,0000\% \\
    \hline
    3 & S0 & 1\,600 & 1\,586 & 0 & 14 & 0,8750\% \\
    3 & S1 & 1\,600 & 1\,173 & 0 & 427 & 26,6875\% \\
    3 & S2 & 1\,600 & 1\,381 & 0 & 219 & 13,6875\% \\
    3 & S3 & 1\,600 & 1\,600 & 0 & 0 & 0\% \\
    3 & S4 & 1\,600 & 1\,368 & 0 & 232 & 14,5000\% \\
    \hline
    4 & S3 & 1\,600 & 1\,600 & 0 & 0 & 0\% \\
    4 & S4 & 1\,600 & 1\,249 & 0 & 351 & 21,9375\% \\
    \hline
  \end{tabular}
\end{table}

Para Esc-1, r2--r9 contienen respectivamente 7\,245, 7\,292, 7\,260,
7\,272, 7\,286, 7\,279, 7\,252 y 7\,154 GET; sus conteos HTTP 500 son
0, 26, 10, 31, 0, 0, 1 y 3. Este censo conserva las ocho corridas centrales,
aunque solo r2, r6 y r7 terminaron con código cero; no convierte las cinco
corridas fallidas en muestras válidas del análisis principal.
En Esc-4 cada estrategia aporta 200 solicitudes por repetición: S3 conserva
200 respuestas 2xx en cada una; S4 registra, en orden r2--r9, 62, 52, 93, 5,
52, 5, 69 y 13 HTTP 500, con las restantes respuestas 2xx. La tasa de S4
durante la caída controlada es 351/1600, es decir, 21,9375\% (21,94\%
redondeado). La ausencia de doble adjudicación no implica ausencia de errores
HTTP ni garantiza aislamiento absoluto.

El censo y el detalle por repetición se reproducen con
\texttt{auditar\_http\_arbiter.py --all-scenarios}, pasando como entrada
\path{experimentos/resultados/arbiter/campaign/raw/}. La integridad de los
archivos fuente se comprueba con el manifiesto existente
\path{experimentos/resultados/arbiter/campaign/SHA256SUMS}.
El workflow \path{.github/workflows/ci-cd.yml} verifica los manifiestos,
contrasta las filas de esta tabla con el censo y las latencias con
\texttt{summary/runs.json}, y ejecuta el análisis offline de las campañas ISO
correctivas versionadas. Una prueba negativa altera una latencia en memoria
y exige que la auditoría la rechace; no modifica el RAW preservado.

S0 es el baseline formal y control negativo; S3 no lo es. La comparación
descriptiva de 26,6875\% frente a 0\% corresponde a S1 frente a S3 en Esc-3.
Respecto del baseline formal S0, la comparación correcta en ese escenario es
26,6875\% frente a 0,8750\%. Esta asimetría impide atribuir superioridad
operacional global a S1--S4.

En Esc-2, Esc-3 y Esc-4, cada \texttt{HTTP\_ERROR} conserva \texttt{sent\_ns}, \texttt{received\_ns} y
\texttt{latency\_ms}. Los HTTP 500 fueron incluidos en los promedios por corrida:
la latencia publicada representa todos los intentos bajo carga, incluidas las
respuestas 500, y no solo adjudicaciones exitosas. No hubo exclusión selectiva
de fallos de la población principal.

En S3, los tiempos de recuperación de las ocho muestras analíticas fueron 0,
0, 1, 0, 0, 0, 0 y 0 ms: media 0,125 ms, mediana 0 ms y rango 0--1 ms. En siete
casos no se observó demora medible con la resolución milimétrica disponible;
esto no equivale a recuperación instantánea. En S4, la recuperación observada
por healthcheck presentó media 3252,11 ms, mediana 3171,85 ms, desviación
estándar 2188,52 ms y rango 1116,75--5949,68 ms. Esta medición demuestra el
retorno del servicio al estado saludable, no el tiempo interno de convergencia
de Raft.

El oráculo evaluó 96 muestras analíticas. I1 (ausencia de doble adjudicación)
produjo 80 PASS y 16 FAIL, todos estos últimos en S0; I2 (exactamente un
adjudicatario) produjo 96 PASS; e I3 (equipo fuera de mantenimiento), 96 PASS.
I4 (liberación exactamente una vez tras cancelar) e I5 (acceso solo con reserva
vigente) quedaron ambos en 96 \texttt{NOT\_OBSERVED}; no se interpretan como
aprobados.

Para la tasa de doble adjudicación, las comparaciones prerregistradas entre S0
como grupo A y cada estrategia S1--S4 como grupo B, tanto en Esc-2 como en
Esc-3, obtuvieron $U=64$, $p=0,0007775304469403844$ y $A_{12}=1,0$, con ocho muestras por
grupo. $A_{12}$ expresa la probabilidad de que A produzca un valor mayor que B:
en esta métrica, un valor mayor es desfavorable, de modo que 1,0 representa el
peor comportamiento de S0 y no su superioridad. Las comparaciones de latencia
mostraron diferencias entre S0 y S3 en Esc-2 ($U=4$, $p=0,00328$,
$A_{12}=0,0625$) y entre S0 y S1--S3 en Esc-3 ($p\leq0,00114$ y
$A_{12}\leq0,015625$); S0 frente a S4 no mostró diferencia estadística en
Esc-2 ($p=0,7527$) ni Esc-3 ($p=0,5286$). Estos contrastes no compensan la
inseguridad observada en S0.

El índice de Jain fue 1,0000, pero no discrimina equidad en este diseño: cada
corrida usa una única franja y un único equipo y, en S1--S4, normalmente existe
un solo adjudicatario. Se trata de un resultado degenerado, no de evidencia de
superioridad. Los rechazos innecesarios medidos fueron cero porque todas las
solicitudes contenciosas clasificadas como \texttt{REJECTED} compartieron equipo
y franja con una adjudicación confirmada. \texttt{HTTP\_ERROR} no se convierte
en confirmación ni en rechazo: por ello, cero dobles adjudicaciones y cero
rechazos innecesarios pueden coexistir con fallos HTTP y no significan que todas
las solicitudes terminaran correctamente. Los valores $U=64$,
$p=0,0007775304469403844$ y $A_{12}=1,0$ siguen siendo matemáticamente
reproducibles y los 500 no invalidan su cálculo, pero tampoco prueban
disponibilidad equivalente. La conclusión queda limitada a seguridad observada
bajo la carga ofrecida, no a corrección o superioridad operacional absoluta.
Además, siete corridas de Esc-1 fallaron ante respuestas HTTP 500 y este
escenario aporta solo $n=3$ muestras válidas; I4 e I5 no
fueron observados; S3 está limitado por la resolución temporal; y S4 mide
healthcheck, no consenso directamente. La evidencia solo permite afirmar que no
se observaron modificaciones en las tablas productivas que fueron verificadas,
no impacto cero absoluto sobre toda la producción.

La evidencia versionada de la campaña está disponible en
\path{experimentos/resultados/arbiter/campaign/}. Corresponde al SHA
experimental \texttt{e43967eda31410b7fef060cfc7903cea94efbe96} e incluye
\path{manifest/checkpoint.json}, \path{summary/runs.json},
\path{summary/runs.csv}, \path{analysis/comparisons.json},
\path{analysis/oracle/}, \path{raw/} y \path{SHA256SUMS}. Este último contiene
315 entradas verificadas.

\subsection{Evaluación ISO/IEC 25010 seleccionada}
\label{subsec:tabla-iso25010}

Observabilidad apoya la eficiencia, la tasa de fallos y el diagnóstico, pero no
certifica ISO 25010 por sí sola. La Tabla~\ref{tab:resultado-iso25010} integra
solo características con evidencia disponible.

\begin{table*}[t]
  \centering
  \small
  \caption{Resultado de las características ISO/IEC 25010 seleccionadas.}
  \label{tab:resultado-iso25010}
  \begin{tabular}{|p{\dimexpr.17\textwidth-2\tabcolsep-1.2\arrayrulewidth\relax}|p{\dimexpr.18\textwidth-2\tabcolsep-1.2\arrayrulewidth\relax}|p{\dimexpr.20\textwidth-2\tabcolsep-1.2\arrayrulewidth\relax}|p{\dimexpr.29\textwidth-2\tabcolsep-1.2\arrayrulewidth\relax}|p{\dimexpr.16\textwidth-2\tabcolsep-1.2\arrayrulewidth\relax}|}
    \hline
    \textbf{Característica} & \textbf{Métrica} & \textbf{Objetivo} & \textbf{Evidencia/resultado} & \textbf{Estado} \\
    \hline
    Eficiencia & p95 nominal, 50 usuarios & IC95 superior $<500$ ms & Campaña poblada; media 9,952465 ms e IC95 [9,595591; 10,309339] ms & CUMPLE en escenario medido \\
    \hline
    Eficiencia secundaria & p99 nominal & IC95 superior $<750$ ms & Campaña poblada; media 19,744969 ms; IC95 [16,871382; 22,618556] ms & CUMPLE en escenario medido \\
    \hline
    Fiabilidad acotada & HTTP 5xx correctivo poblado, 50 usuarios/1 h & IC95 superior $<1\%$ & Media 0,005870\%; IC95 [-0,002938; 0,014677]\% en r2--r9 & CUMPLE en escenario \\
    \hline
    Disponibilidad & Tiempo apto/tiempo total & $\geq99,5\%$ & Diez horas ejecutadas, sin indicador temporal operacionalizado & NO CONCLUYENTE \\
    \hline
    Seguridad & Decisiones dinámicas por Gateway & 100\%; cero falsos permitidos/rechazados y flaky & 7/7 distintos; repetibilidad 3/3; Wilson [0,645670; 1] & CUMPLE en alcance reducido \\
    \hline
    Mantenibilidad & Cobertura por componente & Umbrales previos de líneas y puertas secundarias & Todos cumplen salvo Android: 38,3407\% de líneas frente a 70\% & NO CUMPLE \\
    \hline
    Compatibilidad & Suite Web en tres motores & 100\%; cero fallidos, omitidos y flaky & Chromium, Firefox y WebKit: 8/8; repetibilidad 3/3; Wilson [0,675592; 1] & CUMPLE en suite Web \\
    \hline
  \end{tabular}
\end{table*}

Seguridad, mantenibilidad y compatibilidad se midieron en E3 sobre el SHA
experimental
\texttt{fa7d75ec0f75573938bf46ed6a68f0aee99606ac}. Los commits posteriores
incorporan análisis y documentación; no cambian retrospectivamente el software
medido. El resultado consolidado se conserva en
\path{experimentos/resultados/analisis-e3.json} y su relación con el
protocolo, productores y raw se explicita en
\path{experimentos/resultados/TRAZABILIDAD-E3-E4.md}.

En seguridad se realizaron tres repeticiones de los mismos siete casos. Los 7/7
casos distintos fueron correctos, con proporción 1 e IC95 Wilson
[0,6456695648259365; 1]; 3/3 repeticiones fueron exitosas e idénticas. No se
suman como 21 observaciones independientes. La conclusión queda limitada a
login, reservas y solicitudes (incluidas revisión/propuesta), fixtures, Gateway
y entorno ensayados. La matriz estática de 198 operaciones es antecedente de
superficie, no evidencia de 198 pruebas dinámicas.

En compatibilidad, Chromium, Firefox y WebKit obtuvieron 8/8 casos aprobados
por motor, sin fallidos, omitidos ni flaky; 3/3 resúmenes por motor fueron
idénticos. Cada IC95 Wilson es [0,6755924350132556; 1]. La evidencia acredita
comportamiento cross-browser de la suite Web ensayada, no una evaluación
integral de coexistencia/interoperabilidad ISO ni de Android.

\subsection{Mantenibilidad: evidencia y alcance}
\label{subsec:mantenibilidad}

La \textbf{modularidad} se sustenta en cinco microservicios con propiedad de
datos diferenciada y en clientes Web y Android separados. Dentro de los
servicios, controladores, casos de uso, puertos y adaptadores reducen el impacto
de sustituir persistencia o clientes HTTP. Esta separación no es absoluta: el
dominio Académico todavía expone tipos de paginación de Spring, limitación
descrita en la Sección~\ref{sec:capas-solid}.

La \textbf{analizabilidad} se apoya en Checkstyle, ESLint, logs estructurados,
métricas, trazas y reportes Surefire, JaCoCo y V8. La arquitectura registra
decisiones mediante ADR y emplea Repository, Strategy, State, Observer, Facade,
Factory y Singleton en puntos concretos; estos patrones aportan vocabulario y
puntos de variación, pero su sola presencia no prueba bajo acoplamiento.

La \textbf{modificabilidad} se evidencia mediante contratos Pact y gates por
servicio: un cambio se valida antes de construir imágenes ligadas al SHA. La
\textbf{testabilidad} cuenta con unitarias, integración Spring/Room,
Testcontainers, pruebas de seguridad, Pact, Vitest, Playwright e instrumentadas
Android. Los umbrales son verificables en configuración: 70\% de líneas para
Auth, Usuarios, Académico y Gateway; 80\% de líneas y 48\% de ramas para
Reservas; y 70\% para statements, ramas, funciones y líneas Web.

La Tabla~\ref{tab:mantenibilidad-e3} presenta todas las puertas oficiales E3.
Cada métrica produjo el mismo valor en las tres repeticiones, por lo que
$s=0$ y el IC95 t es degenerado [media; media]. Esto refleja repetición idéntica
del proceso de medición sobre el mismo SHA, no certeza universal ni independencia
estadística entre líneas.

\begin{table*}[ht]
  \centering
  \small
  \caption{Resultados oficiales de mantenibilidad E3, tres repeticiones por métrica.}
  \label{tab:mantenibilidad-e3}
  \begin{tabular}{|l|l|r|r|l|}
    \hline
    \textbf{Componente} & \textbf{Métrica} & \textbf{Media/IC95 (\%)} & \textbf{Umbral (\%)} & \textbf{Decisión} \\
    \hline
    Auth & Líneas & 88,0422 & 70 & CUMPLE \\
    Usuarios & Líneas & 84,0601 & 70 & CUMPLE \\
    Académico & Líneas & 83,1609 & 70 & CUMPLE \\
    Reservas & Líneas & 84,3634 & 80 & CUMPLE \\
    Reservas & Ramas & 56,7256 & 48 & CUMPLE \\
    Gateway & Líneas & 88,8889 & 70 & CUMPLE \\
    Web & Líneas (HISTÓRICA / E3) & 89,91 & 70 & CUMPLE \\
    Web & Ramas (HISTÓRICA / E3) & 73,69 & 70 & CUMPLE \\
    Web & Funciones (HISTÓRICA / E3) & 81,87 & 70 & CUMPLE \\
    Web & Sentencias (HISTÓRICA / E3) & 85,96 & 70 & CUMPLE \\
    Android & Líneas & \textbf{38,3407} & 70 & \textbf{NO CUMPLE} \\
    \hline
  \end{tabular}
\end{table*}

Las cifras Web de la Tabla~\ref{tab:mantenibilidad-e3} corresponden a una
medición E3 anterior que no incluía la totalidad de \texttt{src/} y quedan
marcadas como \textbf{HISTÓRICA / E3}; no deben interpretarse como la cifra
vigente. La medición vigente del alcance completo \texttt{apps/web/src/**/*.\{ts,tsx\}}
se documenta en la Tabla~\ref{tab:mantenibilidad-web-vigente}.

Para Backend, las líneas se recalcularon desde los campos \texttt{LINE\_MISSED}
y \texttt{LINE\_COVERED} de los \texttt{jacoco.csv} canónicos, mediante
\texttt{jacoco\_summary} de \texttt{analizar\_e3.py}, sin depender de XML no
versionados. En
\path{experimentos/resultados/evidencia-e3-canonica/e3_mantenibilidad/},
las tres repeticiones \texttt{rep-01} a \texttt{rep-03} conservan
\texttt{usuarios/report/jacoco.csv} con 1735 líneas cubiertas y 329 no cubiertas:
$100(1735/2064)=84,06007751937985\%$; para
\texttt{reservas/report/jacoco.csv}, 2730 cubiertas y 506 no cubiertas:
$100(2730/3236)=84,36341161928307\%$. CI recalcula ambas coberturas desde
esos CSV y contrasta su redondeo con la tabla. La complejidad
JaCoCo se reporta solo de forma descriptiva (missed+covered): Auth 352/56
clases, media 6,285714 y máximo 34; Usuarios 858/107, 8,018692 y 68;
Académico 1204/141, 8,539007 y 57; Reservas 2029/203, 9,995074 y 178;
Gateway 61/9, 6,777778 y 28. No es una puntuación de calidad. El paquete E3
no conserva un artefacto cuantitativo de Checkstyle con exit code que permita
reconstruir retrospectivamente esa medición; el CI posterior es evidencia separada.

\begin{table*}[ht]
  \centering
  \small
  \caption{\textbf{MEDICIÓN WEB VIGENTE} de cobertura, alcance completo de \texttt{apps/web/src/}.}
  \label{tab:mantenibilidad-web-vigente}
  \begin{tabular}{|l|r|r|l|}
    \hline
    \textbf{Métrica} & \textbf{Resultado (\%)} & \textbf{Umbral (\%)} & \textbf{Decisión} \\
    \hline
    Sentencias (statements) & 81,89 & 70 & CUMPLE \\
    Ramas (branches) & 70,25 & 70 & CUMPLE \\
    Funciones (functions) & 78,43 & 70 & CUMPLE \\
    Líneas (lines) & 84,87 & 70 & CUMPLE \\
    \hline
  \end{tabular}
\end{table*}

Comando exacto: \texttt{npm run test:coverage} (equivalente a \texttt{vitest run --coverage}).
Alcance: \texttt{apps/web/src/**/*.\{ts,tsx\}}. Umbral mínimo del gate: 70\%
para statements, branches, functions y lines (sin reducir). Suite: 43 archivos,
296 pruebas, 3/3 repeticiones consecutivas en PASS. Contexto de rama:
\texttt{fix/eval2-7-web-ivan}, base en el commit \texttt{c10abc0} del repositorio
oficial \texttt{gleiston-guerrero/Entrega-final-del-PFC}.

La decisión global experimental de mantenibilidad es \textbf{NO CUMPLE
únicamente por Android}, cuya cobertura de líneas fue
38,34070796460177\%, inferior al umbral previo de 70\%. Las cifras del SHA
\texttt{e43967eda31410b7fef060cfc7903cea94efbe96}, incluida la referencia
Android de 45,69\%, permanecen como antecedentes históricos de CI y no se usan
como resultado oficial E3.

\subsection{Análisis crítico}
\label{subsec:analisis-critico-final}

Las fortalezas principales son una p95 nominal con margen respecto de
500 ms, una tasa 5xx correctiva poblada cuyo IC95 superior permanece por debajo
de 1\%, instrumentación de las tres señales, evidencia raw trazable, pruebas
multicapa y gates CI/CD ligados al SHA. La
arquitectura permite relacionar un cambio con pruebas, imagen y despliegue, y el
manifiesto hace detectables alteraciones en el conjunto experimental.

Las limitaciones impiden una conclusión total. El p99 histórico de eficiencia presenta
gran dispersión aunque cumple su umbral previo; la rampa contiene un máximo
superior a 14 segundos; Mobile permanece bajo 70\% de cobertura; la campaña histórica de
fiabilidad registró decenas de miles de 401 por repetición; la primera campaña
con refresh quedó limitada por un dataset vacío, mientras la campaña poblada
registró 0 HTTP 401 y 42 HTTP 5xx preservados. La disponibilidad temporal no
fue operacionalizada. A ello se suman la falta de una captura runtime final
del panel Mobile E2E poblado bajo carga, falta de rollback automático,
ausencia de Pact Broker, Testcontainers no uniforme y artefactos de calidad no
persistidos para todos los componentes. Estas restricciones se documentan como
amenazas a la generalización, no se ocultan mediante promedios o extrapolaciones.

\subsection{Reflexión ética}
\label{subsec:reflexion-etica}

SCLI trata perfiles de estudiantes y docentes, asignaciones institucionales,
reservas, incidentes y registros de asistencia. Aunque el sistema no está
integrado con la matrícula académica oficial ni con fuentes externas de la
universidad, estos datos internos permiten inferir presencia, horarios y
actividad. Por ello deben recopilarse únicamente los atributos necesarios para
operar el laboratorio, limitarse su conservación al propósito declarado y
evitar que observabilidad o logs repliquen credenciales, tokens o información
personal innecesaria.

La trazabilidad de aprobaciones y asistencia favorece la rendición de cuentas,
pero también puede convertirse en vigilancia excesiva si se usa para finalidades
distintas de las informadas. El control por roles reduce exposición, sin
eliminar la obligación de revisar periódicamente permisos y accesos. Las
métricas deben agregarse cuando sea posible y no utilizarse para elaborar
perfiles personales de productividad. Estas decisiones siguen los principios
de privacidad, prevención del daño, honestidad y responsabilidad profesional
del Código de Ética de la ACM \cite{acm2018code}.

Las decisiones académicas o disciplinarias no deben automatizarse a partir de
una ausencia, alerta o métrica aislada. SCLI aporta evidencia operacional; una
persona responsable debe verificar contexto, errores de identificación,
indisponibilidad técnica y reglas institucionales antes de actuar. El equipo
mantiene responsabilidad sobre diseño, validación y comunicación de
limitaciones, incluida la ausencia de integración con registros académicos
oficiales.

La responsabilidad profesional también alcanza la seguridad del producto y la
veracidad de su evaluación. Los permisos deben aplicarse en Backend y no solo
ocultar controles en las interfaces; los incidentes y trazas deben conservarse
sin exponer datos personales; y ninguna automatización debe desplazar el juicio
humano en decisiones que afecten a estudiantes o docentes. En coherencia con el
Código de Ética y Práctica Profesional de Ingeniería de Software, aprobado
conjuntamente por ACM e IEEE Computer Society \cite{acmieee1999secode}, el equipo
debe comunicar tanto los resultados favorables como los fallos y límites. Por
ello se mantienen las siete corridas fallidas de ARBITER, los invariantes no
observados, el análisis histórico p99 incorrectamente poblado y los 401 de la
campaña de fiabilidad, en
lugar de convertir evidencia desfavorable en cumplimiento.

\subsection{Declaración de uso de inteligencia artificial generativa}
\label{subsec:declaracion-ia}

Durante el desarrollo y la redacción se empleó inteligencia artificial
generativa como herramienta de apoyo. El código y el texto resultantes fueron
revisados por integrantes del equipo y contrastados con el repositorio, las
pruebas y la evidencia disponible. Una salida generada por IA no se considera
por sí sola evidencia experimental: los resultados reportados proceden de los
archivos de protocolo, datos crudos, resúmenes y verificaciones versionadas. La
selección técnica, la interpretación de resultados y la responsabilidad final
corresponden al equipo FUVV; no se atribuye a la herramienta autoría ni decisión
autónoma.

\subsection{Amenazas a la validez}
\label{subsec:amenazas-validez}

\subsubsection{Validez interna}

\textbf{Amenaza:} calentamiento de JVM, cachés y recursos compartidos entre
contenedores pueden alterar la latencia. \textbf{Efecto:} diferencias entre
repeticiones que no provienen del cambio evaluado. \textbf{Mitigación aplicada:}
protocolo previo, diez repeticiones nominales, conservación de r1 y r10 y uso
uniforme de r2--r9 para el análisis central. \textbf{Riesgo residual:} no se
aislaron por completo CPU, red, cachés ni pausas del runtime en el host.

\textbf{Amenaza:} las preguntas PI1 y PI2 se formalizaron en el manuscrito
después de ejecutadas las campañas, aunque el protocolo, la población y los
umbrales de decisión ya estaban versionados previamente.
\textbf{Efecto:} existe riesgo de sesgo retrospectivo al expresar preguntas de
investigación con conocimiento de resultados ya observados.
\textbf{Mitigación aplicada:} se conserva la cronología Git de los criterios
previos, no se repitieron las campañas para favorecer las respuestas y las
reglas permiten resultados negativos.
\textbf{Riesgo residual:} las preguntas no fueron prerregistradas literalmente
antes de la ejecución; por tanto, sus respuestas deben interpretarse como una
formalización posterior aplicada a criterios experimentales preexistentes.

\textbf{Amenaza:} la disponibilidad HTTP fue asimétrica entre estrategias
ARBITER; Esc-3/S1 alcanzó 26,6875\% de respuestas 500, mientras Esc-3/S3 tuvo
0\% y el baseline S0 0,8750\%. \textbf{Efecto:} una estrategia puede evitar
dobles adjudicaciones porque controla la concurrencia, pero también porque parte
de la carga falla antes de confirmar. \textbf{Mitigación aplicada:} se conservan
y publican todos los intentos y códigos HTTP. \textbf{Riesgo residual:} la tasa
de doble adjudicación no operacionaliza disponibilidad ni corrección total.

\textbf{Amenaza:} la tasa 5xx no operacionaliza tiempo disponible.
\textbf{Efecto:} el cumplimiento del umbral de errores podría interpretarse
erróneamente como disponibilidad. \textbf{Mitigación aplicada:} la campaña
correctiva poblada completó diez ventanas válidas de una hora con 0 HTTP 401 y
42 HTTP 5xx. La campaña histórica con 401 y la primera correctiva con dataset
vacío y 0 HTTP 5xx permanecen publicadas por separado; no se calcula
disponibilidad a partir de ninguna de las tres.
\textbf{Riesgo residual:} el tiempo apto frente al tiempo total sigue siendo
desconocido.

\subsubsection{Validez de constructo}

\textbf{Amenaza:} la campaña histórica de PI1 usó un dataset vacío y no
ejercitó \texttt{GET /api/v1/reservas/\{id\}}.
\textbf{Efecto:} medir únicamente listados vacíos no representa todas las
consultas autenticadas de solo lectura declaradas por la pregunta.
\textbf{Mitigación aplicada:} la campaña poblada corrige esa degeneración
para el cierre post-evaluación, con ambas rutas y 20 reservas controladas;
el antecedente histórico y su alta dispersión se conservan por separado.
\textbf{Riesgo residual:} no se reconstruyó exactamente el dataset histórico
y no se afirma comparabilidad estricta entre poblaciones. Las 20 reservas
no implican representatividad universal; las conclusiones de PI1 y PI2
quedan acotadas al escenario, datos, carga y entorno medidos.

\subsubsection{Validez externa}

\textbf{Amenaza:} la evidencia procede de una VM y un conjunto de datos y carga
controlados, no de una instalación distribuida representativa de todas las
condiciones institucionales. \textbf{Efecto:} p95, p99 y tasa de errores pueden
cambiar con otra red, hardware, volumen o distribución de usuarios.
\textbf{Mitigación aplicada:} se documentan versión, SHA, herramientas,
parámetros y datos crudos para repetir el escenario. \textbf{Riesgo residual:}
los resultados solo se generalizan al escenario medido.

\textbf{Amenaza:} el número de escenarios de carga y compatibilidad es limitado.
\textbf{Efecto:} pueden quedar combinaciones de navegador, dispositivo o
contención sin observar. \textbf{Mitigación aplicada:} se usan tres motores Web,
un emulador Android y una rampa exploratoria además del caso nominal.
\textbf{Riesgo residual:} dispositivos reales y redes adversas permanecen sin
cubrir; aunque ARBITER S0--S4 fue ejecutado, Esc-1 conservó solo tres muestras
analíticas válidas y el dataset experimental no representa toda la demanda de
un campus real.

\subsubsection{Validez de constructo}

\textbf{Amenaza:} Locust, Prometheus interno y la métrica móvil miden latencia
en puntos diferentes. \textbf{Efecto:} tratarlas como equivalentes mezclaría
latencia desde generador, latencia interna del backend y tiempo HTTP observado
desde Android; además, ninguna de ellas mide renderizado de UI.
\textbf{Mitigación aplicada:} se reportan por separado y se explica el punto de
medición de cada una. \textbf{Riesgo residual:} falta una captura runtime final
del dashboard con la métrica móvil poblada bajo carga en la VM.

\textbf{Amenaza:} el estimando de latencia ARBITER mezcla respuestas válidas y
HTTP 500, con tasas de fallo diferentes entre estrategias. \textbf{Efecto:} no
puede interpretarse como latencia exclusiva de adjudicaciones exitosas ni como
comparación de disponibilidad equivalente. \textbf{Mitigación aplicada:} la
población principal conserva todos los intentos y se informa el censo HTTP por
separado. \textbf{Riesgo residual:} restringir ahora el análisis principal a
2xx sería una redefinición post hoc no prerregistrada. Un eventual análisis
solo-2xx sería exclusivamente una sensibilidad secundaria y no sustituiría el
resultado principal.

\textbf{Amenaza:} cobertura, número de pruebas y configuración del pipeline no
equivalen por sí solos a calidad completa. \textbf{Efecto:} declarar
cumplimiento por conteo puede omitir seguridad, concurrencia o comportamiento
semántico. \textbf{Mitigación aplicada:} se combinan pruebas unitarias,
contratos, integración, E2E y evidencia experimental, y las cifras variables se
remiten al reporte del SHA. \textbf{Riesgo residual:} Seguridad,
Compatibilidad y Mantenibilidad permanecen parciales donde falta evidencia
integral actual. Esas auditorías históricas parciales no sustituyen las
decisiones experimentales E3 reportadas en la
Sección~\ref{subsec:tabla-iso25010}. La trazabilidad contractual incorporada
posteriormente se presenta en la Sección~\ref{sec:contratos-openapi}.

\subsection{Trazabilidad final de la Entrega 4}
\label{subsec:trazabilidad-final}

\begin{table}[ht]
  \centering
  \scriptsize
  \caption{Matriz final de trazabilidad respecto de la rúbrica.}
  \label{tab:trazabilidad-final}
  \begin{tabular}{|p{1.1cm}|p{3.0cm}|p{4.1cm}|p{2.2cm}|p{1.7cm}|}
    \hline
    \textbf{Req.} & \textbf{Implementación} & \textbf{Evidencia} & \textbf{Sección} & \textbf{Estado} \\
    \hline
    D1.1 & Capas y puertos & Servicios; dependencia Spring Page/Pageable en dominio Académico & Arquitectura/SOLID & PARCIAL \\
    \hline
    D1.2 & Siete patrones reales & ADR-005 y clases Repository, Factory, Strategy, State, Observer, Facade, Singleton & Patrones & CUMPLE \\
    \hline
    D2.1 & SPA Web funcional y logout remoto & React, rutas, reservas, JWT y revocación & Aplicación Web & CUMPLE \\
    \hline
    D2.2 & Calidad Web & Umbrales V8 al 70\%; Playwright en tres motores & Pruebas Web & Resultado por SHA en CI \\
    \hline
    D3.1 & Cliente Android & Compose/MVVM, Room, QR; incidentes locales y FCM no E2E & Aplicación Mobile & PARCIAL \\
    \hline
    D3.2 & Calidad Mobile & Antecedente del SHA \texttt{e43967e}: 45,69\%; E3 oficial: 38,3407\%, bajo 70\% & Pruebas Mobile & NO CUMPLE en E3 \\
    \hline
    D4.1 & Contratos Pact & Consumidores y verificadores versionados; resultado por SHA & Pact & EVIDENCIADO EN CI \\
    \hline
    D4.2 & Integración CockroachDB & Testcontainers, Flyway, cifrado e idempotencia & Testcontainers & CUMPLE \\
    \hline
    D5.1 & Pirámide de pruebas & Backend, Web, Mobile, integración y E2E verdes en \texttt{e43967e}; Mobile conserva medición documental bajo 70\% & Estrategia de pruebas & PARCIAL \\
    \hline
    D5.2 & CI/CD con gates y SHA & Workflows, GHCR, APK, deployment; sin rollback y Actions por tag & CI/CD & PARCIAL \\
    \hline
    D6.1 & Métricas, trazas y logs & Prometheus, OTel, Tempo, Loki, cAdvisor y Grafana & Observabilidad & CUMPLE \\
    \hline
    D6.2 & Dashboard de seis paneles & JSON canónico; panel 6 con latencia HTTP E2E móvil reportada desde Android; falta captura runtime poblada bajo carga & Dashboard & PARCIAL \\
    \hline
    D7.1 & Experimento ISO 25010 & Eficiencia 10/10; fiabilidad correctiva 10/10 y E2 CUMPLE $<1\%$; disponibilidad temporal no operacionalizada & ISO/experimentos & PARCIAL \\
    \hline
    D7.2 & Experimento ARBITER & 130 planificadas; 123 completadas; 99 muestras analíticas válidas & ARBITER/oráculo & PARCIAL \\
    \hline
  \end{tabular}
\end{table}

\subsection{Conclusiones}
\label{subsec:conclusiones-e4}

La Entrega 4 consolidó un sistema distribuido con cinco aplicaciones Backend,
clientes Web y Android, persistencia CockroachDB y controles de seguridad
integrados. La arquitectura por capas, los patrones aplicados y los contratos
HTTP permiten separar responsabilidades, aunque permanece deuda puntual en el
dominio Académico y en capacidades Mobile sin backend remoto.

La calidad cuenta con evidencia reproducible: pruebas unitarias, integración
real con CockroachDB, Pact, Playwright, Android instrumentado, Locust y un
pipeline que condiciona imágenes y despliegue. Backend y Web tienen gates de
cobertura configurados; la evidencia Mobile del HEAD permanece bajo 70\%, por
lo que no se afirma cumplimiento global ni se atribuyen al HEAD cifras de un
reporte anterior.

La eficiencia nominal poblada cumple los criterios p95 y p99 y no registró HTTP 5xx en
las muestras centrales; la rampa muestra picos de cola. La campaña correctiva
poblada de fiabilidad cumple el criterio acotado de 5xx: r2--r9 presenta media
0,005870\,\% e IC95 [-0,002938; 0,014677]\,\%. Las diez repeticiones
pobladas preservan 0 HTTP 401 y 42 HTTP 5xx; los 401 de la campaña histórica y
la campaña intermedia con dataset vacío permanecen como antecedentes. Este resultado no demuestra
disponibilidad de 99,5\,\%; la conclusión profesional sigue siendo parcial
para la evaluación ISO completa.

El experimento ARBITER confirmó la utilidad del control negativo: S0 presentó
dobles adjudicaciones en las 16 muestras analíticas de contención. No se
observaron dobles adjudicaciones con S1--S4 en las corridas ejecutadas, sin que
esto constituya una garantía universal ni superioridad operacional: los fallos
HTTP fueron asimétricos y alcanzaron 26,6875\% en Esc-3/S1. El resultado expresa
seguridad observada bajo la carga ofrecida y no disponibilidad equivalente. De
130 corridas planificadas se
completaron 123 y quedaron 99 muestras analíticas válidas; I4 e I5 no fueron
observados y las mediciones de recuperación conservan las limitaciones de
resolución y healthcheck descritas en la Sección~\ref{subsec:arbiter-resultados}.

La observabilidad mejora la capacidad de diagnóstico mediante Prometheus,
OpenTelemetry, Tempo, Loki, cAdvisor y Grafana. Su dashboard satisface la
estructura de seis paneles y el panel Mobile ya usa una métrica E2E medida desde
Android y expuesta por Reservas a Prometheus. La evidencia visual final de esa
serie poblada bajo carga en la VM sigue pendiente. En conjunto, la entrega está
sustentada para informe y defensa siempre que las limitaciones experimentales y
técnicas permanezcan explícitas.

\subsection{Trabajo futuro}
\label{subsec:trabajo-futuro}

Las prioridades son elevar la cobertura Mobile con casos de ramas y UI de valor,
diagnosticar mediante una campaña independiente el patrón de respuestas 401,
operacionalizar y validar disponibilidad, investigar y reducir la variabilidad
p99, y capturar bajo carga la latencia E2E móvil ya medida desde Android.
También conviene automatizar rollback, incorporar Pact
Broker con \emph{can-i-deploy}, extender Testcontainers a Académico, persistir
reportes Backend/Web/Playwright y fijar versiones inmutables y uniformes de las
herramientas CI. Estas acciones cierran brechas existentes sin introducir
tecnología ajena a los riesgos observados.
